% Page budget: 1.35 pages | Owns: gantry coordinates, first-principles equations, scheduling, LPV-LFR realisation, and physical parameterisation.
% WRITING GUIDE
% Purpose: Define the physical coordinates and derive the exact continuous-time LPV-LFR baseline used by every later section.
% Start here: docs/Thesis-documentation/LPV_LFR_Derivation/main.tex.
% Supporting material: LPV_LFR_internal_loop_Well_posedness, LPV_Mass_invertible, LPV_LFR_Justification_Drenth, and docs/fp-model-structure.md.
% Mandatory papers: Read Drenth's rational LPV-LFR paper and the LPV-LFR definitions used by Hoekstra before adopting terminology or well-posedness claims.
% Research-plan anchor: Preserve Aspect 1's explicit position dependence; update the original discretisation wording to the final method.
% Current decisions: Continuous-time physics, endogenous scheduler Y, fixed-step RK4, and identifiable parameter combinations.
% Choices to justify: Coordinate frame, Y as scheduler, continuous-time formulation, RK4 step and substeps, full six-channel LFR realization, and reduced physical parameterization. Give the physical or numerical reason for each choice, not only its definition.
% Implementation audit: Read lpv_lfr_baseline/core/{physics,lfr_matrices,lfr_forward}.py, blocks/lfr_param_block.py, and scripts/validate_lfr.py before finalising equations or claims.
% Reference check: Verify the original gantry model, LPV definitions, and LFR framework against primary sources; cite the gantry-specific algebra as this project's own derivation.
% Cautions: The final pipeline uses the full six-channel scheduling loop, not the experimental six-to-four SVD reduction. Resolve the supervisor comment about dual-gantry terminology; do not copy unverified bibliography entries.
% Planned figures: Physical coordinate schematic and compact interconnection between G and Delta(Y).
% Section is finished when: The LFR collapses to the original model, well-posedness assumptions are stated, and notation matches Section 03.
%
% LPV-LFR DERIVATION ROADMAP
% 1. Define q=[X,Theta,Y]^T, the force transformation, measured outputs, assumptions, and the admissible range of Y.
% 2. State M(Y)qdd+Cqd+Kq=f and expose the structural dependence M(Y)=M0+Y M1+Y^2 M2.
% 3. Construct the LFR directly from the polynomial mass equation:
%       M0 qdd + M1 w1 + M2 w2 = f_net,
%       z1=qdd, z2=w1, w1=Y z1, w2=Y z2.
% 4. Collect w=Delta(Y)z with Delta(Y)=Y I6 and give the constant block-level G realization with all signal dimensions.
% 5. Eliminate w and z in a short exactness check to recover the original equation of motion.
% 6. State the well-posedness condition over the complete scheduling range and relate it to invertibility of M(Y); place the longer proof outside the main text.
% 7. Only after the conceptual construction, mention that the implementation
%    evaluates the resulting rational expression in Horner form. Put
%    M(Y)^(-1)=[N0+Y N1+Y^2 N2]/d(Y), the explicit Ni entries, and the
%    determinant coefficients in the appendix or technical supplement.
% 8. Close the derivation by stating that the continuous-time vector field is evaluated first and RK4 supplies the sampled transition used in training.
%
% DERIVATION WARNING
% Do not start from v=M(Y)^(-1)f_net. Roland's direct feedback identifies that as a collapsed inverse interpretation, not the standard LFR construction. Prefer the z1,z2,w1,w2 construction above even where older project notes or the inclusion guide use the inverse-first route. Use LPV/Feedback Supervisor/Roland_Toth_Feedback.md as the governing project note and verify terminology against the primary literature.
% Main text: assumptions, polynomial dependence, direct LFR construction, compact G, exactness, well-posedness statement, and one implementation sentence.
% Appendix/supplement: expanded M_i and N_i matrices, determinant polynomial, longer invertibility/well-posedness proof, and symbolic verification.
\section{Dual-Drive Gantry System and Physics Baseline}

\subsection{System coordinates and equations of motion}
The dual-drive H-type gantry comprises two parallel actuators that carry a
cross-arm through flexible joints and a third actuator that moves the payload
along the arm (Fig.~\ref{fig:gantry_lfr}(a)) \cite{garcia2013model}. Let
$q=[X\;\Theta\;Y]^\top$ denote the generalized coordinates and
$\qact=[X_1\;X_2\;Y]^\top$ the actuator coordinates. They satisfy
\begin{equation}
  \qact = P^\top q,
  \qquad
  P = \begin{bmatrix} 1 & 1 & 0 \\ L_b/2 & -L_b/2 & 0 \\ 0 & 0 & 1 \end{bmatrix},
  \label{eq:Ptransform}
\end{equation}
and virtual work gives the generalized force $u=P\uact$. This coordinate choice
avoids a closed kinematic-chain description and exposes the coupling caused by
the position-dependent load distribution \cite{garcia2013model}. The measured
output is $\qact=P^\top q$.

The adopted lumped-parameter model is
\begin{equation}
  M(Y)\ddot{q}(t) + C\dot{q}(t) + K q(t) = P\uact(t),
  \label{eq:eom}
\end{equation}
The payload position enters the inertia matrix without derivatives and does so
quadratically,
\begin{equation}
  M(Y) = M_0 + Y M_1 + Y^2 M_2 ,
  \label{eq:Msplit}
\end{equation}
where $C$ and $K$ are constant. Only the yaw coordinate has a restoring
stiffness; $X$ and $Y$ are rigid-body coordinates. The expanded matrices and
the parameter values inherited from the supplied reference implementation are
reported in the appendix; these parameters were not re-identified in this work.

For positive component masses and inertias, $M(Y)\succ0$ for all $Y\in\R$; a
proof by Sylvester's criterion \todo{Check that we actually use Sylvester's criterion...} is given in the appendix. Hence, $M(Y)$ is
invertible throughout the operating range, as required below.

The baseline deliberately excludes the effects to be learned by the augmentation.
In the simulated plant these are a payload-mounted absorber and, for the friction
benchmark, Coulomb friction. Garc\'ia-Herreros et al. likewise identify unmodelled
supporting-structure vibration, motor detent forces, and friction variation along
the stroke \cite{garcia2013model}; these effects are not equated with the synthetic
absorber used here. \todo{Determine and justify the operating range of $Y$; distinguish
the mechanical stroke, the simulated-data range, and the range covered by the measured
ILC records.}

\begin{figure*}[t]
  \centering
  \begin{subfigure}{0.65\textwidth}
    \centering
    % ===========================================================================
% Final thesis figure: lumped-parameter model of the dual-drive H-type gantry
% with the simulated payload absorber. Geometry adapted from
% Garcia-Herreros et al. (2013), Fig. 2.
%
% Requires in the preamble:
%   \usepackage{tikz}
%   \usetikzlibrary{arrows.meta,calc,decorations.pathmorphing}
%
% Conventions
%   * Coordinates are in drawing units; the x/y keys of the tikzpicture set the
%     unit, so the drawing scales while the labels keep the document's own font
%     and size.  At 0.3mm the drawing is 1:1 with Garcia Fig. 2 and the figure
%     is about 275pt wide, i.e. half of a two column text width (figure*).
%   * The origin is the cross arm centre of mass (the blue coordinate X).
%   * Parts on the cross arm are drawn in its rotated frame: u along the beam
%     (positive to the right), v perpendicular to it.  The measured coordinate
%     Y points LEFT in this view, so Y = -u, and the absorber, which sits at
%     +L0 in Y, is drawn to the LEFT of the payload.
%   * Changing \th (yaw) or \Yp (payload position) moves every attached part,
%     including its labels.  With \Yp > 0 the absorber ends up between the CoM
%     and the payload, and the F_Y label then needs a manual nudge.
%   * Schematic, not to scale: the absorber offset L0 and the mass sizes are
%     drawn for legibility.  The absorber rides on the payload, so it may
%     overhang the beam end.
%
% \gantrydelta: 1 adds an arrow for the absorber displacement relative to the
% payload.
% ===========================================================================
\providecommand{\gantrydelta}{1}
\begin{tikzpicture}[
    x=0.3mm, y=0.3mm,                        % original Garcia scale (1 unit
                                             % = 1pt of Garcia Fig. 2)
    line width=0.5pt, line cap=round, line join=round,
    >={Triangle[length=3.4pt,width=2.9pt]},
    dim/.style={line width=0.4pt, dash pattern=on 3pt off 1.1pt on 0.4pt off 1.1pt},
    ext/.style={line width=0.35pt},
    cl/.style={gray!55, line width=0.25pt,
               dash pattern=on 4pt off 1pt on 0.3pt off 1pt},
    ref/.style={blue!70, line width=0.3pt, dash pattern=on 1.2pt off 0.8pt},
    abs/.style={draw=black, fill=black!15, line width=0.5pt}, % absorber: shaded
    absl/.style={font=\small},                                % k_a, c_a labels
  ]
  % ================ parameters ================
  \def\th{9.25}      % yaw angle Theta [deg]
  \def\Lb{152.5}     % cross arm length
  \def\Wb{24.5}      % cross arm width
  \def\Yp{-18}       % payload position along the cross arm (u); Y = -u
  \def\Lh{26}        % payload width
  \def\Hh{24}        % payload height
  \def\rw{3.7}       % payload wheel radius
  \def\Wa{20}        % absorber width
  \def\Ha{20}        % absorber height
  \def\ga{30}        % spring and damper gap between absorber and payload
  \def\Mw{34.7}      % actuator mass box size
  \pgfmathsetmacro\yE{0.5*\Lb*sin(\th)}       % height of the beam ends
  \pgfmathsetmacro\hb{0.5*\Wb}
  \pgfmathsetmacro\vh{\hb+2*\rw+0.5*\Hh}      % payload centre height (v)
  \pgfmathsetmacro\ua{\Yp-0.5*\Lh-\ga-0.5*\Wa}  % absorber centre (u)
  \pgfmathsetmacro\vY{\vh+0.5*\Hh+30}          % height of the Y dimension line
  \pgfmathsetmacro\xM{0.5*\Lb*cos(\th)+22.3+0.5*\Mw}  % |x| of actuator centres
  % Lane for the d dimension and the F_Y arrow: on the CoM side of the payload
  % and clear of the red Y reference line, so it follows the sign of \Yp.
  \pgfmathsetmacro\ud{ifthenelse(\Yp<0,9,-9)}
  \pgfmathsetmacro\sd{ifthenelse(\Yp<0,1,-1)}

  % ===================== cross arm =====================
  \begin{scope}[rotate=\th]
    \draw (-\Lb/2,-\hb) rectangle (\Lb/2,\hb);
    \draw[cl] (-0.57*\Lb,0) -- (0.58*\Lb,0);
    \coordinate (E1) at ( \Lb/2,0);
    \coordinate (E2) at (-\Lb/2,0);

    % ==== payload m_h on its wheels ====
    \draw[fill=white] (\Yp-\Lh/2,\vh-\Hh/2) rectangle (\Yp+\Lh/2,\vh+\Hh/2);
    \draw (\Yp-7,\hb+\rw) circle (\rw)  (\Yp+7,\hb+\rw) circle (\rw);
    \node at (\Yp,\vh) {$m_h$};

    % ==== payload absorber: mass m_a, spring k_a, damper c_a ====
    \pgfmathsetmacro\xa{\ua+0.5*\Wa}          % absorber side attachment
    \pgfmathsetmacro\xh{\Yp-0.5*\Lh}          % payload side attachment
    \draw[abs] (\ua-\Wa/2,\vh-\Ha/2) rectangle (\ua+\Wa/2,\vh+\Ha/2);
    \node at (\ua,\vh) {$m_a$};
    \pgfmathsetmacro\xc{0.5*(\xa+\xh)}        % centre of the gap
    % spring k_a on top
    \draw[decorate, decoration={zigzag, segment length=4, amplitude=3.5,
          pre length=5, post length=5}] (\xa,\vh+6) -- (\xh,\vh+6);
    \node[absl, anchor=south, inner sep=1.5pt] at (\xc,\vh+10) {$k_a$};
    % damper c_a below: rod, filled piston, open cylinder, rod
    \draw (\xa,\vh-6) -- (\xc+1,\vh-6);
    \fill (\xc-1.4,\vh-8.6) rectangle (\xc+0.8,\vh-3.4);
    \draw (\xc-4,\vh-2) -- (\xc+4,\vh-2) -- (\xc+4,\vh-10) -- (\xc-4,\vh-10);
    \draw (\xc+4,\vh-6) -- (\xh,\vh-6);
    \node[absl, anchor=north, inner sep=0.5pt]        % tucked up under the
          at (\xc,\vh-10.9) {$c_a$};                   % cylinder, clear of the beam
    \ifnum\gantrydelta>0                      % displacement relative to m_h
      \draw[->] (\ua,\vh+\Ha/2+4.5) -- ++(-14,0);
      \node[anchor=south east, inner sep=1pt]
            at (\ua-3,\vh+\Ha/2+5.5) {$\delta_a$};
    \fi
    % ==== yaw reference at the payload: axis perpendicular to the beam ====
    \coordinate (P) at (\Yp,-1.5);
    \draw[ref] (P) -- ++(0,\vY+34);

    % ==== payload force F_Y ====
    \draw[->] ({\ud+\sd*11},\vh+4) -- ({\Yp+\sd*\Lh/2},\vh+4);
    \node[inner sep=2pt] at ({\ud+\sd*20},\vh+4) {$F_Y$};

    % ==== d : centreline to payload centre of mass ====
    \draw[<->] (\ud,0) -- (\ud,\vh);
    \node[inner sep=2pt] at ({\ud-\sd*4.5},0.78*\vh) {$d$};

    % ==== Y : CoM to payload (red) ====
    \begin{scope}[red]
      \draw[dim] (0,\hb+2) -- (0,\vY+6);   % the blue payload axis is the
                                            % left extension of this dimension
      \draw[->] (0,\vY) -- (\Yp,\vY);
      \node[anchor=south] at ({0.5*\Yp},\vY+1) {$Y$};
    \end{scope}

    % ==== flexible joints and L_b ====
    \draw[dim] ( \Lb/2,-\hb-2) -- ++(0,-24)   (-\Lb/2,-\hb-2) -- ++(0,-24);
    \draw[<->] (-\Lb/2,-\hb-18) -- node[pos=0.72, below=1pt] {$L_b$}
               ( \Lb/2,-\hb-18);
    \node[anchor=north west, inner sep=2.5pt] at ( \Lb/2,-\hb-19) {$k_{b1},\,c_{b1}$};
    \node[anchor=north east, inner sep=2.5pt] at (-\Lb/2,-\hb-19) {$k_{b2},\,c_{b2}$};
  \end{scope}

  % ============ Theta at the payload: vertical reference and arc ============
  \draw[ref] (P) -- ++(0,{(\vY+34)*cos(\th)});
  \draw[ref, line width=0.4pt]
        ($(P)+(90:\vY+22)$) arc[start angle=90, end angle=90+\th, radius=\vY+22];
  \node[blue, anchor=east, inner sep=2pt]
        at ($(P)+(90+\th:\vY+24)$) {$\Theta$};

  % ===================== generalized coordinates X, Theta (blue) =============
  \begin{scope}[blue]
    \draw[->] (0,-5.2) -- (0,-26);
    \node[anchor=west, inner sep=2pt] at (1.5,-21) {$X$};
    \draw (4.6,0) -- (36,0);
    \draw (32,0) arc[start angle=0, end angle=\th, radius=32];
    \node[anchor=west, inner sep=2.5pt] at (37,{36*sin(0.5*\th)}) {$\Theta$};
  \end{scope}
  % centre of mass m_b
  \fill[white] (0,0) circle (4.3);
  \fill (0,0) -- (4.3,0) arc[start angle=0, end angle=-90, radius=4.3] -- cycle
        (0,0) -- (-4.3,0) arc[start angle=180, end angle=90, radius=4.3] -- cycle;
  \draw (0,0) circle (4.3);
  \node[anchor=north east, inner sep=1pt] at (-3.4,-3.4) {$m_b$};

  % ===================== actuators, guides, forces =====================
  % \s = +1 : actuator 1 (right) ; \s = -1 : actuator 2 (left)
  \foreach \s/\i in {1/1,-1/2}{
    \pgfmathsetmacro\y{\s*\yE}                 % row height (symmetric)
    \coordinate (M\i) at (\s*\xM,\y);
    \draw (\s*\xM-\Mw/2,\y-\Mw/2) rectangle (\s*\xM+\Mw/2,\y+\Mw/2);
    \node at (M\i) {$m_\i$};

    % flexible joint: spiral spring between beam end and actuator
    \pgfmathsetmacro\xs{\s*(0.5*\Lb*cos(\th)+11)}
    \draw (\xs,\y) -- (E\i);
    \draw[line width=0.4pt] plot[domain=0:540, samples=100, smooth]  % 1.5 turns
          ({\xs+(3.4+3.9*\x/540)*cos(90*(1-\s)-\s*(\x-540))},
           {\y +(3.4+3.9*\x/540)*sin(90*(1-\s)-\s*(\x-540))});
    \draw (\xs+\s*7.3,\y) -- (\s*\xM-\s*\Mw/2,\y);

    % linear guide: two rollers against a hatched wall
    \pgfmathsetmacro\xb{\s*(\xM+\Mw/2+5.35)}
    \draw (\xb,\y+11.6) circle (5.35)  (\xb,\y-11.6) circle (5.35);
    \draw (\xb+\s*5.6,\y-20) -- (\xb+\s*5.6,\y+20);
    \foreach \k in {0,...,5}
      \draw[ext] (\xb+\s*5.6,\y+18-7.5*\k) -- ++(\s*6.5,-6.5);

    % actuator force
    \draw[->] (\s*\xM,\y+\Mw/2+30) -- (\s*\xM,\y+\Mw/2);
    \node[anchor=south] at (\s*\xM,\y+\Mw/2+31) {$F_{X\i}$};

    % measured position X_i (red)
    \pgfmathsetmacro\xr{\s*(\xM+\Mw/2+23.5)}  % clear of the wall hatching
    \draw[red] (\xr,\y+4.5) -- (\xr,\y-4.5)  (\xr,\y) -- (\xr+\s*5,\y);
    \draw[red,->] (\xr+\s*5,\y) -- (\xr+\s*5,\y-26);
    \node[red, anchor=north] at (\xr+\s*5,\y-27) {$X_\i$};
  }

  % ===================== inertial frame (top view) =====================
  \begin{scope}[shift={(72,-58)}]
    \draw[->] (0,-3.2) -- (0,-19);          % x points down
    \draw[->] (-3.2,0) -- (-19,0);          % y points left
    \draw[fill=white] (0,0) circle (3.2);   % z into the page: circle with a cross
    \draw[ext] (-2.26,-2.26) -- (2.26,2.26)  (-2.26,2.26) -- (2.26,-2.26);
    \node[anchor=east, inner sep=2pt]  at (-19,0.3) {$y$};
    \node[anchor=west, inner sep=2pt]  at (3.6,0.8) {$z$};
    \node[anchor=north, inner sep=2pt] at (0,-20)  {$x$};
  \end{scope}
\end{tikzpicture}

    \caption{Lumped-parameter gantry and coordinate definitions.}
    \label{fig:gantry_system}
  \end{subfigure}%
  \hfill%
  \begin{subfigure}{0.33\textwidth}
    \centering
    % LFR interconnection of the baseline: constant G closed by the scheduling block.
% Adapted from LPV/LFR-derivation-supervisor.tex, relabelled for continuous time
% and for the gantry scheduling variable Y.
% Requires: \usepackage{tikz} \usetikzlibrary{arrows.meta}
\begin{tikzpicture}[
    >={Latex[length=2.2mm]},
    line width=0.5pt,
    block/.style={draw, rectangle, minimum width=15mm, minimum height=11mm},
    lbl/.style={font=\footnotesize, inner sep=1.5pt},
  ]
  \node[block] (G)     at (0,0)    {$G$};
  \node[block] (D)     at (0,17mm) {$\Dsch(Y)$};

  \pgfmathsetmacro{\po}{2}     % port offset from the block centre [mm]
  \pgfmathsetmacro{\rx}{16}    % routing distance from the centre [mm]

  % scheduling input
  \draw[->] (0,34mm) -- (D.north);
  \node[lbl, left] at (0,31mm) {$Y$};

  % w : Delta -> G  (left branch)
  \draw[->] (D.west) -- (-\rx mm,17mm) -- (-\rx mm,{\po mm}) -- (-7.5mm,{\po mm});
  \node[lbl, left] at (-\rx mm,10mm) {$w$};

  % z : G -> Delta  (right branch)
  \draw[->] (7.5mm,{\po mm}) -- (\rx mm,{\po mm}) -- (\rx mm,17mm) -- (D.east);
  \node[lbl, right] at (\rx mm,10mm) {$z$};

  % external input and output
  \draw[->] (-25mm,{-\po mm}) -- (-7.5mm,{-\po mm});
  \node[lbl, left] at (-25mm,{-\po mm}) {$u$};
  \draw[->] (7.5mm,{-\po mm}) -- (28mm,{-\po mm});
  \node[lbl, right] at (28mm,{-\po mm}) {$q$};
\end{tikzpicture}

    \caption{Self-scheduled baseline interconnection.}
    \label{fig:lfr}
  \end{subfigure}
  \caption{Physics baseline. (a) Lumped-parameter schematic adapted from the
  mechanical arrangement in \cite{garcia2013model}; the absorber subsystem is present
  only in the simulated data-generating plant. (b) The constant interconnection
  $G$ is closed by $\Dsch(Y)=YI_6$, where $Y$ is a component of the state.}
  \label{fig:gantry_lfr}
\end{figure*}

\subsection{Self-scheduled LPV-LFR representation}\label{sec:lfr}
The payload position is both a measured coordinate and a model state. Since $Y$
is a system state rather than an exogenous signal, the resulting formulation is
quasi-LPV \cite{toth2010modeling}. It is self-scheduled through the known map
$p=Y$. The state-space matrices depend rationally on $Y$ through $M(Y)^{-1}$;
an LFR represents this dependence exactly using a constant interconnection and
an explicit scheduling block \cite{drenth2025thesis}.

Collecting the net generalized force acting on the inertia,
$\fnet = u-C\dot{q}-Kq$, and substituting \eqref{eq:Msplit}
into \eqref{eq:eom} gives
\begin{equation}
  M_0\,\ddot{q} + M_1 w_1 + M_2 w_2 = \fnet,
  \qquad
  w_1 = Y z_1, \quad w_2 = Y z_2,
  \label{eq:latent}
\end{equation}
with the loop signals $z_1=\ddot{q}$ and $z_2=w_1$. Stacking
$z=[z_1^\top\;z_2^\top]^\top$ and $w=[w_1^\top\;w_2^\top]^\top$, the scheduling
variable acts only through
\begin{equation}
  w = \Dsch(Y)\,z, \qquad \Dsch(Y) = Y I_6 ,
  \label{eq:delta}
\end{equation}
while all remaining relations are independent of $Y$. The direct six-channel
realization is retained because its two three-dimensional channels correspond
directly to the linear and quadratic terms in \eqref{eq:Msplit}. The baseline is
the lower linear fractional transformation $\mathcal{F}_\ell(G,\Dsch(Y))$ of
\begin{equation}
  \begin{bmatrix} \dot{x} \\ z \\ q \end{bmatrix}
  =
  \underbrace{\begin{bmatrix}
      A   & B_w    & B_u \\
      C_z & D_{zw} & D_{zu} \\
      C_q & 0      & 0
  \end{bmatrix}}_{G}
  \begin{bmatrix} x \\ w \\ u \end{bmatrix},
  \label{eq:lfr}
\end{equation}
where $x=[q^\top\;\dot{q}^\top]^\top$, $n_x=6$, $n_z=n_w=6$, and
$n_u=n_q=3$. The matrices in $G$ depend on the physical parameters but not on
$Y$; their expressions in terms of $M_0^{-1}$, $M_1$, $M_2$, $C$, and $K$ are
listed in the appendix. The transformations $u=P\uact$ and
$\qact=P^\top q$ connect this internal realization to the actuator interface.
Substituting
\eqref{eq:delta} back into \eqref{eq:latent} returns
$(M_0+YM_1+Y^2M_2)\ddot{q}=\fnet$, so the realization reproduces \eqref{eq:eom}
exactly.

The interconnection is well posed when $I-D_{zw}\Dsch(Y)$ is nonsingular for
every admissible $Y$ \cite{drenth2025thesis}. With $M_0$ invertible, this condition
is equivalent to nonsingularity of $M(Y)$. The positive-definiteness result above
therefore establishes well-posedness throughout the operating range.

The continuous-time mechanical model is retained instead of pre-discretizing a
set of frozen models. A fixed-step RK4 scheme evaluates the state-dependent
scheduler at its intermediate stages and supplies the sampled transition used in
training and evaluation. The sampling interval and number of integration
substeps are specified in Section~\ref{sec:setup}. The implementation eliminates
the algebraic loop analytically and evaluates the resulting rational expression
in Horner form; the coefficient matrices and determinant polynomial are given in
the appendix.

\subsection{Physical parameters and identifiable combinations}
\begin{itemize}
\item Distinguish the fourteen raw parameters from the ten combinations identifiable through the gantry equations
\item Optimise reduced positive parameter combinations and map them back to interpretable masses, stiffnesses, damping, geometry, and centre-of-mass terms
\item Show the Jacobian from raw parameters to identifiable combinations; naive parameter merging is not a valid recovery argument
\item Use exact-model parameter recovery only as an internal implementation/initialization check; do not allocate a main result to it
\item Use relative combination error as the interpretability metric and disclose the special normalisation required when a combination crosses zero
\item Should state that Lh is part of the transformation (P-transform) so we dont identify this one.
\end{itemize}
\todo{Should the raw-to-reduced Jacobian be printed in the appendix or shown only as rank and singular values?}
